% MAIN SNAPSHOT LSPV-5 1012 2026-10-02T10:12:00-07:00
\documentclass{article}\usepackage{amsmath,amssymb}\begin{document}
\section*{LSPV-5 MAIN 1012}
\textbf{STATUS: OPEN.} P versus NP remains open. No unrestricted SAT lower bound, arbitrary-algorithm transfer, nonrelativizing hinge, or nonalgebrizing hinge is established.

\section*{CORPUS RECONCILIATION}
Surveyed 285 repository paths; root main; AGENTS/CODEX/skill governance; TASK-001--007; MAIN through 0954; current LSPV-1--4 material; BITRANK/MATRIXCERT/EXCESS/TURANEX; HEC/VSC-HEC/G-PPC-HEC/CRS-PPC/CCSPC/RCCN/Search-G; TRACECLOSE1; PHCLOSE1/PRIMEPH1; prime provenance; experiments/results; governance; GROK/CLAUDE material.

New ACR1 result: authenticated causal replay does not generically escape existing closures. Deterministic polynomial replay plus polynomial authentication is P-closed by TRACECLOSE1. Polynomial-length fixed-PH replay validity is selectable in FP under $P=NP$ by PH-PROV-CLOSE1. ACRTREE1 gives a finite hostile control showing explicit adaptive strategy-tree size can be exponential while a complete response strategy has a succinct polynomial representation.

Literature reconciliation: incrementally verifiable computation shows composable verification can remain efficient; proofs of sequential work rely on cryptographic/sequentiality assumptions; pebbling, online communication and pointer chasing give model-specific lower bounds; $AP=PSPACE$ identifies the polynomial-alternation boundary but does not force an arbitrary static SAT decider to materialize an explicit strategy.

No stale claim is strengthened. REACQ1/TRACECLOSE1/PH-PROV-CLOSE1, COVERBREAK1, EXCESS1, BITRANK characteristic-transfer failure, SGLAUD1 failure and BOOL hostile controls remain retained. Ordinary switching, generic canonical binding, generic provenance mass, raw incidence/tape geometry and cryptographic hardness transfers remain retired.

\section*{Shortest surviving proof chain}
\[
\boxed{\text{exact source-typed SU/vMS family}}
\to\boxed{\operatorname{rank}_{\mathbb F_2}\Phi=\omega(m)\text{ or classified exceptions}}
\to\boxed{\text{source-bound adaptive SAME\_REALIZATION necessity}}
\]
\[
\to\boxed{\text{charged retention/reacquisition surviving replay and succinct strategies}}
\to\boxed{\text{model-matched reconstruction/static compiler}}
\to\boxed{\text{nonrelativizing/nonalgebrizing hinge}}
\to\boxed{\text{unconditional contradiction}}.
\]
Every load-bearing boxed arrow remains OPEN or CONDITIONAL.

\section*{ACR status}
The only surviving ACR candidate commits one exact source realization $R$ before adaptive legal SU/vMS restrictions and requires later openings from that same $R$. A useful theorem must prove a representation-independent retention/reacquisition lower bound surviving static reread, succinct strategy representation, recomputation, memoization/pointer jumping, semantic-oracle replacement, algebraic-oracle replacement, INTERP and RECOMP. Polynomial alternation outside PH is a boundary, not a lower bound.

\section*{ONE SINGLE TEAM TARGET}
For one exact reconstruction-compatible SU/vMS polynomial-bit orbit, build the source-transcribed binary operator $\Phi$ and legal slice matrices, then instantiate a source-bound adaptive-replay game on those same legal restrictions. Jointly measure characteristic-two rank, response-signature diversity, and minimum retained/reacquired information. Attack with succinct strategies, full static reread, recomputation, maximal semantic oracle and algebraic oracle. ACR remains CANDIDATE unless it survives these attacks.

P versus NP remains OPEN.
\end{document}
% END_OF_VERIFIED_MAIN::1012
